\chapter{lamina}\label{ch:lamina}
\begin{goals}
\textbf{In this chapter} you learn what lamina is for, what it reads, the two things it never does
(change your file, use the network), and the one idea behind how it handles size.
\end{goals}

A block model arrives as a 3~GB CSV. A spreadsheet stops at a million rows; a text editor opens it and
stops responding; a database wants it imported first. lamina opens it in about the time it takes to read
the file once, shows you its rows, and lets you filter, sort, derive columns and report on it, without
loading it.

That is the whole purpose: to look at a large table you have been given, find out what is in it, check
that it is what it claims to be, and take out the part you need. lamina is a \term{viewer}. It reads; it
doesn't edit.

\section{What it reads}\label{sec:reads}

\begin{description}
\item[Delimited text] CSV, TSV and any delimited \file{.txt} or \file{.dat}, at any size: comma, tab,
semicolon, pipe or runs of whitespace; a header row or none; quoted fields; a block of \code{\#}
comment lines before the header; GSLIB and Geo-EAS files with their column names in a preamble; a
decimal comma; UTF-8, Windows-1252, Latin-1 or UTF-16 (\cref{ch:open}).
\item[Datamine \file{.dm}] block models and other Datamine tables, read directly at any size, with no
conversion step (\cref{sec:dm}).
\item[Apache Parquet] the columnar format pandas, Polars, DuckDB and Spark write, typed from the file's
own schema and read by range (\cref{sec:parquet}).
\item[Excel workbooks] \file{.xlsx} and \file{.xlsm}, one sheet at a time, typed as Excel typed them
(\cref{sec:xlsx}).
\item[Archives] \file{.zip}, \file{.tar}, \file{.gz}, \file{.zst}, \file{.bz2} and \file{.xz}: lamina
looks inside and opens an entry without unpacking the archive to disk (\cref{sec:archives}).
\item[Anything else] plain text opens as a table of lines; a binary file is named as binary rather than
shown as noise.
\end{description}

\section{It never changes your file}\label{sec:read-only}

lamina opens files read-only. Nothing you do in it writes to the file you opened: a filter, a sort, a
hidden column, a calculated column, a number format, a colour scale. All of these belong to the
\term{view}, the way you are looking at the table, and they last as long as the window does (or as long
as a lens keeps them, \cref{ch:lens}).

When you want a file out, \menu{Data>Export…} writes a new one, somewhere you choose
(\cref{ch:export}). Opening a Datamine \file{.dm} and exporting it is how you convert it to CSV, at any
size.

\section{It never touches the network}\label{sec:never-network}

lamina is \term{Sealed}. The page carries a browser security policy that forbids every network
connection, so it cannot send anything anywhere, and it doesn't fetch anything either: no fonts, no
scripts, no update check, no telemetry. Your data stays in the page, on your machine. The policy is
part of the file, so you can read it, and the build checks it before every release (\cref{ch:sealed}).

\section{How size stops mattering}\label{sec:windowed}

lamina never holds the whole file. When you open one it reads it once, from start to end, and keeps a
small index: where every few thousand rows begin. After that, to show you rows 1{,}204{,}000 to
1{,}204{,}040, it looks up the nearest index entry, reads that stretch of the file, and decodes only
the rows on screen. Scrolling, jumping to a row and resizing all work this way. A filter, a sort or a
statistic is one more pass through the file, with only what the pass needs kept in memory: a list of the
matching rows, the sort keys, the running sums.

So memory stays roughly flat as files grow, and the first look at a file costs one read of it. The
second time you open the same file, the index comes from a cache and the table is there at once. \Cref{ch:huge} goes through the details and the costs.

\begin{why}
Windowing is the design, not an optimisation. A tool that loads the file first works until the file is
bigger than the memory you have, and resource-estimation files routinely are. Reading by windows has a
cost (a sort needs a pass, and some operations have limits, \cref{app:limits}), but it means there is no
file size at which lamina stops working.
\end{why}

\section{What it is not}

lamina is not a spreadsheet: cells can't be edited and there are no formulas in cells, only calculated
columns derived from a whole row. It is not a database: it doesn't join tables or keep anything once you
close it. It is not a 3D viewer: micro, its sibling, is the tool for seeing a block model in space. And it is not a geostatistics package: it reports and
checks, it doesn't estimate.

What it is, is the first thing to open when a file arrives. Within a minute you know how many rows and
columns it has, what each column holds, which values are missing, whether the codes kept their leading
zeros, how the grades are distributed, and how much tonnage sits above a cutoff.
